
\begin{abstract}
\blfootnote{\vspace{-1.5em}\\
Corresponding author: Heng Guo\\
\email{guoheng@bupt.edu.cn} \vspace{0.1cm}\\ 
$1$ \, Beijing University of Posts and Telecommunications, China\\
$2$ \, National Institute of Informatics, Japan\\
$3$ \, vivo BlueImage Lab, vivo Mobile Communication Co., Ltd., Shanghai, China\\
$4$ \, State Key Laboratory for Multimedia Information Processing, School of Computer Science, Peking University, China\\
$5$ \, National Engineering Research Center of Visual Technology, School of Computer Science, Peking University, China
}
	Accurate reconstruction of reflective surfaces remains a fundamental challenge in computer vision, with broad applications in real-time virtual reality and digital content creation. Although 3D Gaussian Splatting (3DGS) enables efficient novel-view rendering with explicit representations, its performance on reflective surfaces still lags behind implicit neural methods, especially in recovering fine geometry and surface normals. To address this gap, we propose \polgs, a physically-guided polarimetric Gaussian Splatting framework for fast reflective surface reconstruction. Specifically, we integrate a polarized BRDF (pBRDF) model into 3DGS to explicitly decouple diffuse and specular components, providing physically grounded reflectance modeling and stronger geometric cues for reflective surface recovery. 
	Furthermore, we introduce a \emph{depth-guided visibility mask acquisition} mechanism that enables angle-of-polarization (AoP)-based tangent-space consistency constraints in Gaussian Splatting without costly ray-tracing intersections. 
	This physically guided design improves reconstruction quality and efficiency, requiring only 	about 10 minutes of training. Extensive experiments on both synthetic and real-world datasets validate the 	effectiveness of our method.
Code page: \href{https://github.com/PRIS-CV/PolGS\_plus}{https://github.com/PRIS-CV/PolGS\_plus}.
\keywords{Polarization-based vision; Reflective surface reconstruction; 3D Gaussian Splatting}
\end{abstract}
\vspace{-2em}
